% Final: full independent proof/source audits by /root and /root/frobenius.
\paragraph{C.3ab (Proved): integral top-cycle finiteness for surface cones with an explicit vanishing hypothesis.}
Let $k$ be an algebraically closed field of characteristic zero,
$V$ a smooth projective integral surface over $k$, and $L$ a very
ample line bundle whose full section ring
\[
 S=\bigoplus_{j\ge0}H^0(V,L^j)
\]
is normal and generated in degree one. Put
\[
 R=(S_{S_+})^{\mathrm{sh}},\qquad
 X=\operatorname{Spec}R\setminus\{\mathfrak m_R\},\qquad
 g=\gcd\{\deg(c_1(L)\cdot C):C\subset V
                         \text{ an integral curve}\}.
\]
Assume the following cohomological condition, with \emph{absolute}
K\"ahler differentials over $\mathbb Q$:
\begin{equation}\label{eq:C3ab-absolute-vanishing}
 H^2(V,\Omega^1_{V/\mathbb Q}\otimes L^j)=0
                         \qquad(j\ge1).
\end{equation}
Then the natural maps from C.3e are isomorphisms, integrally:
\begin{equation}\label{eq:C3ab-answer}
 \mathrm{CH}^2(X)\simeq
 \frac{\mathrm{CH}_0(V)}{c_1(L)\mathrm{CH}^1(V)},
 \qquad
 B^2(X)\simeq\mathbb Z/g\mathbb Z.
\end{equation}
Here $B^2$ uses ordinary algebraic equivalence on actual fixed
products with parameter curves. The integer $g$ is positive.
Condition \eqref{eq:C3ab-absolute-vanishing} holds for every such
surface over $\overline{\mathbb Q}$. Over an arbitrary $k$ as above,
it holds if $H^2(V,L^j)=0$ for all $j\ge1$; in particular this
includes $K3$ surfaces and surfaces with $p_g=0$.
This statement does not settle arbitrary three-dimensional
strict local germs, arbitrary complex surface cones, or the top
groups in local dimension at least four.

\begin{proof}
\emph{1. The resolution and its actual infinitesimal neighbourhoods.}
The degree-one hypothesis identifies the blow-up of
$\operatorname{Spec}S$ at its vertex with
$W=\operatorname{Tot}(L^{-1})$, as proved by charts in C.3e.
Set
\[
 Z=W\times_{\operatorname{Spec}S}\operatorname{Spec}R,
 \qquad i:V\hookrightarrow Z,
 \qquad p:Z\longrightarrow V.
\]
Then $Z$ is a regular integral three-dimensional scheme, proper
over $R$, and $Z\setminus V=X$. The divisor $V$ is Cartier,
its normal bundle is $L^{-1}$, $p\circ i=1$, and
$\mathfrak m_R\mathcal O_Z=\mathcal O_Z(-V)$.
The morphism $p$ is the limit of the smooth relative-dimension-one
maps in C.3e; in particular it is flat and has the indicated
flat pullback on cycles. For $a\ge1$ write
\[
 V_a=Z\times_R R/\mathfrak m_R^a.
\]
Strict henselization does not change the quotients of the local
cone by powers of its maximal ideal, because the residue field
is already algebraically closed. Thus these are the split
thickenings
\begin{equation}\label{eq:C3ab-split-thickenings}
 V_a=\operatorname{Spec}_V
       (\mathcal O_V\oplus L\oplus\cdots\oplus L^{a-1}),
\end{equation}
with the truncated graded multiplication. Their projections to
$V$ split restriction on $K$-theory and on Picard groups.

\emph{2. Pro cdh descent gives detection at a finite thickening.}
We use nonconnective algebraic $K$-theory, with $K_0$ the
Grothendieck group of perfect complexes. Apply
\citet[Theorem 3.7(i), author pp.~21--22]{D_Morrow2016}
to the abstract blow-up square formed by $Z$, $\operatorname{Spec}R$
and their closed fibres. All schemes are noetherian,
quasi-excellent $\mathbb Q$-schemes of finite Krull dimension.
The resulting exact sequence of pro abelian groups includes
\[
 K_0(R)\longrightarrow
 \{K_0(R/\mathfrak m_R^a)\}_a\oplus K_0(Z)
 \longrightarrow\{K_0(V_a)\}_a.
\]
The first groups are all canonically $\mathbb Z$, and the first
arrow is $n\mapsto(n,n[\mathcal O_Z])$. Quotienting by this graph
shows that
\begin{equation}\label{eq:C3ab-pro-monomorphism}
 K_0(Z)\longrightarrow\{K_0(V_a)\}_a
\end{equation}
is a monomorphism in the abelian category of pro abelian groups.
This assertion is stronger than merely separating a fixed
individual element in an inverse limit. Indeed, the kernel pro
group is represented by the decreasing subgroups
$\ker(K_0(Z)\to K_0(V_a))$, whose transition maps are inclusions.
A pro-zero system of this form has a zero term. Consequently
there is $a_0$ such that
\begin{equation}\label{eq:C3ab-finite-detection}
 K_0(Z)\longrightarrow K_0(V_a)\quad\text{is injective for }a\ge a_0.
\end{equation}
No assertion that arbitrary $K$-theory commutes with a formal
inverse limit is used.

\emph{3. The relative $K_2$ calculation retains integral information.}
Let $\mathcal K_q(V_a,V)$ denote the Zariski sheaf of relative
$K_q$ on their common underlying space $V$. Nilpotence gives
$\mathcal K_0(V_a,V)=0$. Determinant and the finite logarithm give
\[
 \mathcal K_1(V_a,V)=1+\mathcal I_a
       \simeq\bigoplus_{j=1}^{a-1}L^j,
 \qquad \mathcal I_a=\ker(\mathcal O_{V_a}\to\mathcal O_V).
\]
For a nilpotent ideal in a $\mathbb Q$-algebra, all relative
$K$-groups are uniquely divisible
\citep[Consequence 1.4, p.~392]{D_Weibel1982}.
Therefore the rational comparison of
\citet[Main Theorem, p.~348]{D_Goodwillie1986}
is an integral isomorphism of the underlying abelian groups here.
In degree two it identifies the relative sheaf with relative
$HC_1$ over $\mathbb Q$.

The latter calculation can be made explicitly. On a trivializing
affine open of $V$, write $A[t]/(t^a)$ for
\eqref{eq:C3ab-split-thickenings}. For a commutative
$\mathbb Q$-algebra $B$, the degree-one cyclic complex gives
$HC_1(B)=\Omega^1_{B/\mathbb Q}/dB$: this follows directly from
the cyclic relation in degree one, or from the low-degree
Connes exact sequence $B\xrightarrow{d}\Omega^1_{B/\mathbb Q}
\to HC_1(B)\to0$.
The positive homogeneous part of degree $j$, $1\le j<a$, is
\[
 \frac{t^j\Omega^1_{A/\mathbb Q}\oplus
                 A t^{j-1}dt}
      {\{t^j\,db+j b t^{j-1}dt:b\in A\}}.
\]
Every element has a unique horizontal representative:
\[
 [t^j\omega+b t^{j-1}dt]
          \longmapsto t^j(\omega-j^{-1}db).
\]
Equivalently, inclusion of the horizontal summand is an
isomorphism onto this quotient. The relation
$a t^{a-1}dt=0$ removes degree $a$, and there are no further
positive degrees. Under a change of fibre coordinate the
horizontal inclusion transforms by the $j$th power of the
line-bundle transition function. Thus these local computations
give the global isomorphism
\begin{equation}\label{eq:C3ab-relative-K2}
 \mathcal K_2(V_a,V)\simeq
       \bigoplus_{j=1}^{a-1}
                      \Omega^1_{V/\mathbb Q}\otimes L^j.
\end{equation}
In particular no extra conormal summand has been retained after
quotienting by exact differentials.

\emph{4. The only relative $K_0$ classes are determinant classes.}
Apply Zariski descent to the fibre of $K(V_a)\to K(V)$.
Theorem 10.3 of \citet[pp.~383--386]{D_ThomasonTrobaugh1990}
applies to both spectra and hence to their fibre; its spectral
sequence has
\[
 E_2^{s,t}=H^s(V,\mathcal K_t(V_a,V))
                       \Longrightarrow K_{t-s}(V_a,V).
\]
The underlying space has dimension two. On the diagonal
$t=s$ the only possibly nonzero terms are
$H^1(\mathcal K_1)$ and $H^2(\mathcal K_2)$.
The second vanishes by
\eqref{eq:C3ab-absolute-vanishing} and
\eqref{eq:C3ab-relative-K2}. The first has no possible incoming
or outgoing differential. Naturality for determinant therefore
gives an isomorphism
\[
 K_0(V_a,V)\xrightarrow{\det}\operatorname{Pic}(V_a,V)
         =H^1(V,1+\mathcal I_a).
\]
All these relative groups inject into their absolute groups
because the thickening has a projection to $V$.
It follows that restriction induces an isomorphism
\[
 SK_0(V_a)\xrightarrow{\sim}SK_0(V),\qquad
 SK_0(T)=\ker\bigl(K_0(T)\xrightarrow{(\operatorname{rk},\det)}
                         \mathbb Z\oplus\operatorname{Pic}(T)\bigr).
\]
Combining injectivity in \eqref{eq:C3ab-finite-detection}
with the section $p^*$ gives
\begin{equation}\label{eq:C3ab-SK0}
                 i^*:SK_0(Z)\xrightarrow{\sim}SK_0(V).
\end{equation}

\emph{5. Zero-cycles on $Z$ vanish by the henselian property.}
Every closed point of $Z$ lies on its closed fibre $V$, since
$Z$ is proper over the local ring $R$. Let $z$ be such a point.
In the regular local ring $\mathcal O_{Z,z}$ choose two parameters
whose common prime avoids a local equation for $V$, and close
the resulting integral one-dimensional germ in $Z$. This gives
an integral closed curve $D\subset Z$ through $z$, not contained
in $V$. To justify the parameter choice, complete a local
equation of the smooth Cartier divisor $V$ to a regular system
of three parameters and use the other two.

The closed fibre of $D\to\operatorname{Spec}R$ is
zero-dimensional. Properness and upper semicontinuity of fibre
dimension imply that every fibre is zero-dimensional: otherwise
the closed image of the non-quasi-finite locus would contain
the closed point of the local base. Thus $D$ is finite over
$R$. Its normalization $\widetilde D$ is also finite over $R$
by excellence. A finite algebra over a henselian local ring is
a product of local algebras; because $\widetilde D$ is integral,
it is local. It is consequently the spectrum of a DVR, with
residue field $k$ and unique closed point over $z$.
A uniformizer has divisor equal to that point. Proper
pushforward of this principal divisor proves $[z]=0$ in
$\mathrm{CH}_0(Z)$. Hence
\begin{equation}\label{eq:C3ab-CH0-resolution}
                          \mathrm{CH}_0(Z)=0.
\end{equation}

\emph{6. Codimension two needs no denominator in Riemann--Roch.}
For a smooth variety $T$, let $F^rK_0(T)$ be the filtration by
codimension of coherent support. The rank and determinant
identifications give $F^2K_0(T)=SK_0(T)$, and
\citet[Example 15.3.6, p.~298]{frC4Fulton1998}
gives
\[
 \mathrm{CH}^2(T)\xrightarrow{[D]\mapsto[\mathcal O_D]}
                 F^2K_0(T)/F^3K_0(T)
\]
as an integral isomorphism, inverse to $-c_2$. The two composites
have multiplier $(-1)^{2-1}(2-1)!=-1$, so rationalization is
unnecessary.

These identities also hold on the particular $Z$ used here.
Indeed, $Z$ is the filtered inverse limit of the smooth
threefolds $W_T$ over pointed affine \emph{\'etale}
neighbourhoods of the cone vertex constructed in C.3e.
The transition maps are flat. Vector bundles, their exact
sequences, cycles, and the finite rational-equivalence data
all descend to one such neighbourhood. The support filtrations,
rank, determinant, Chern classes and the displayed isomorphism
therefore pass to this limit. This supplies the precise scope
needed here without invoking a Riemann--Roch theorem for
arbitrary nonalgebraizable regular schemes.

Since $Z$ has dimension three, $F^3K_0(Z)$ is generated by
classes of its closed points. It vanishes by
\eqref{eq:C3ab-CH0-resolution}. Since $V$ is a surface,
$F^3K_0(V)=0$ as well. Compatibility of Chern classes with
pullback and \eqref{eq:C3ab-SK0} now give
\begin{equation}\label{eq:C3ab-CH2-resolution}
 i^!: \mathrm{CH}^2(Z)\xrightarrow{\sim}\mathrm{CH}_0(V),
                      \qquad (i^!)^{-1}=p^*.
\end{equation}

\emph{7. Localization gives the actual ordinary quotient.}
Chow localization gives
\[
 \mathrm{CH}^1(V)\xrightarrow{i_*}\mathrm{CH}^2(Z)
       \longrightarrow\mathrm{CH}^2(X)\longrightarrow0.
\]
Under \eqref{eq:C3ab-CH2-resolution}, self-intersection
identifies the first arrow with multiplication by
$c_1(N_{V/Z})=-c_1(L)$. This proves the first isomorphism
in \eqref{eq:C3ab-answer}.

Flat pullback by $p$, followed by restriction to $X$, sends
ordinary algebraically trivial cycles on $V$ to ordinary
algebraically trivial cycles on $X$: it pulls back each actual
curve-parameter correspondence on the fixed product. Thus the
first isomorphism induces a surjection
\[
 B^2(V)/c_1(L)B^1(V)\longrightarrow B^2(X).
\]
The integral injection already proved in C.3e is this same map,
so it is an isomorphism. On a smooth projective integral
variety over an algebraically closed field, zero-cycles modulo
ordinary algebraic equivalence are classified by degree.
Consequently the last quotient is $\mathbb Z/g\mathbb Z$.
Ampleness makes every degree in the defining gcd positive,
and any hyperplane curve supplies one such degree; hence
$g>0$ and the group is finite.

\emph{8. The field restriction and examples are explicit.}
Akizuki--Kodaira--Nakano vanishing gives
\[
 H^2(V,\Omega^1_{V/k}\otimes L^j)=0\qquad(j\ge1)
\]
because $2+1>\dim V=2$
\citep[Theorem 1.3(a), p.~4; Corollary 6.4 and its proof,
 pp.~55--56]{D_EsnaultViehweg1992}.
If $k=\overline{\mathbb Q}$, absolute and relative differentials
agree. In general smoothness supplies the exact sequence
\[
 0\longrightarrow\Omega^1_{k/\mathbb Q}\otimes_k\mathcal O_V
 \longrightarrow\Omega^1_{V/\mathbb Q}
 \longrightarrow\Omega^1_{V/k}\longrightarrow0.
\]
Thus $H^2(V,L^j)=0$ is a sufficient additional condition.
Cohomology here commutes with the direct sum defining the
possibly infinite-dimensional constant $k$-vector space
$\Omega^1_{k/\mathbb Q}$, as follows from a finite affine
\v Cech cover. For a $K3$ surface this vanishing follows
from $K_V=\mathcal O_V$ and Serre duality, since $L^{-j}$
has no sections. If $p_g=0$, multiplication by a nonzero
section of $L^j$ injects
$H^0(V,K_V\otimes L^{-j})$ into $H^0(V,K_V)=0$,
and Serre duality again applies. The absolute-differential
hypothesis has therefore been verified in the claimed cases,
without replacing it silently by a relative one.

For example, $V=\mathbb P^2_k$ and $L=\mathcal O(d)$ give
$B^2(X)=\mathbb Z/d\mathbb Z$ for every $d\ge1$.
For the Fermat quartic surface in C.3f, choose
$\zeta^4=\eta^4=-1$. The equations
$x_0=\zeta x_1$ and $x_2=\eta x_3$ define a line on the
surface. Its hyperplane degree is one, so $g=1$ and
$B^2(X)=0$ for that particular complex strict local cone.
This is compatible with the non-torsion class of C.3f in
$\mathrm{CH}^2(X)$: its image is ordinarily algebraically
trivial. The two equivalence relations are kept distinct.
\end{proof}
